\section{Discussion}
\label{sec:discussion}

\paragraph{Execution discipline is the lever.} The core claim of this paper is that exploratory tabular QA needs \textbf{execution discipline and auditable runtime control}, not only better retrieval. The companion SANA framework provides direct evidence for the ``not only retrieval'' part: idealizing search and supplying gold sources improves accuracy but does not close the gap to a solved task. The residual gap is what MOMO targets.

\paragraph{The context firewall is the most transferable idea.} Among MOMO's mechanisms, the context firewall is the one we expect to generalize most cleanly to settings outside tabular data lakes. Any agent task that produces large per-step artifacts --- error logs, intermediate scratch outputs, long file reads, exhaustive enumeration --- is a candidate for the same treatment: keep the artifact inside a bounded worker and return only a compact structured summary to the orchestrating context. The artifact, when kept in the working context, does not just consume tokens; it actively degrades the planner's reasoning by anchoring it on details that are no longer relevant.

\paragraph{The partial-success protocol generalizes beyond tabular QA.} The combination of \texttt{required\_outputs} + \texttt{missing\_outputs} + \texttt{status="partial"} is a useful shape for any task where a sub-step produces \emph{enumerable} outputs. Multi-year time series, multi-entity lookups, multi-document summarization, and many information-extraction settings have this structure. The protocol forces the worker to be explicit about what it did \emph{not} produce, which gives the planner a structured signal to act on. Without the protocol, the same partial result either masquerades as success or forces the planner to inspect raw traces to find the gap.

\paragraph{Delegation vs.\ sprint: same goal, different lever.} Both flags address the stopping-rule problem, but they do so at different levels of the agent loop. Sprint forces a single agent to commit to a source and reflect on its commitment at fixed cadences or transitions; the commitment is visible in the agent's own prompt via the persistent \texttt{\#\# CURRENT SPRINT} section. Delegation lifts the commitment boundary up one level entirely: the planner is structurally prevented from doing low-level wrangling, so commitment becomes a property of the contract a worker receives rather than a self-discipline the planner must maintain.

The two have different cost profiles. Sprint adds a small per-cycle prompt overhead and a few mandatory tool calls. Delegation adds the spin-up cost of fresh agents and additional tokens for contracts and returns. We expect delegation to dominate sprint when the task involves multiple distinct source families (the cost of the extra agents is amortized over substantial per-contract work) and to be comparable or slightly worse when the task fits comfortably inside a single source's budget (one inspect contract barely earns back its setup cost). Validating this trade-off empirically is one of the goals of the planned result run.

\paragraph{Why guards belong in the runtime, not in the prompt.} A natural alternative to \texttt{\_FileReferenceGuard} and \texttt{\_InspectSourceGuard} is to add the same rules to the system prompt --- ``always use the exact file path,'' ``do not peek datasets outside the assigned family.'' We have implemented and discarded this alternative. Prompt-level rules are routinely violated under context pressure, and each violation produces a wasted turn plus a polluting traceback. Runtime guards convert the same rule into an immediate, structured contract violation: the worker receives a precise error and can correct on the next turn, with no traceback bleeding into the planner's context. This is the design principle: \emph{push execution discipline as far down the stack as possible}, away from the model and into the runtime, so that compliance does not depend on the model's attention budget.

\paragraph{Implementation as ablation surface.} The choice to implement MOMO as a feature-flagged extension over the baseline \texttt{DataLakeAgent} pays off as an ablation surface. We can run \texttt{results}, \texttt{cot}, \texttt{sprint}, and \texttt{delegation} independently, in pairs, or in combinations (subject to the \texttt{sprint}/\texttt{delegation} mutual-exclusion), and the baseline path is always available as a reference. This is consequential for reproducibility: a reviewer or follow-on researcher can re-run any specific configuration without committing to MOMO as a whole.
